\section{Introduction}
\label{sec: intro}
Vision-centric autonomous driving is a promising direction due to its economic advantages and scalability~\cite{bevformer, tpvformer, surroundocc, occ3d, uniad, vad}. 
Conventional methods typically follow a modular pipeline of perception~\cite{bevdet, bevfusion, gaussianformer, gaussianformer2, quadricformer}, prediction~\cite{occworld, occllama, gaussianworld, doe-1}, and planning ~\cite{uniad, vad, genad, vadv2, gaussianad}, which heavily relies on expensive 3D annotations and thus faces limitations in real-world applications. 
Recently, vision-language-action (VLA) models have emerged to leverage rich world knowledge from large language models to map visual inputs to driving actions~\cite{omnidrive, emma, openemma, orion}, demonstrating remarkable generalization across driving scenarios.

Despite their good generalization across driving scenarios, most existing VLA models only use languages for scene descriptions or decision reasoning and lack flexible instruction-following capabilities~\cite{opendrivevla, orion, autovla, drivemoe, recogdrive}.
They are either trained to imitate an averaged expert policy, or are confined to a closed set of simple navigational commands like ``turn left'' or ``go straight'', failing to generalize to open-ended and flexible natural language instructions.
In contrast, a general driving agent should not only navigate autonomously but also comprehend and execute diverse, user-specified natural language instructions. 
For instance, a user in a hurry might instruct the vehicle to ``overtake the front car to catch the next green light'' rather than adhere to the conservative policy learned from the training data. 

To facilitate the shift from imitation driving to instructional driving, we construct a large-scale driving dataset, \textbf{InstructScene}, with around 100,000 instruction-annotated scenes and the corresponding trajectories built on NAVSIM~\cite{navsim-v1}. 
While a direct way is to train a VLA model on our driving dataset containing rich instructions, we find that it struggles to generate feasible trajectories and follow instructions accurately.
We think this is due to the significant information disparity between the high-dimensional visual-instruction inputs and the low-dimensional action prediction, making it difficult for the model to learn a generalizable mapping from high-level instructions to low-level actions in complex and dynamic environments. 

To address this, we propose a unified \textbf{vision-language-world-action} model, \model, for joint instruction-based generation and planning.
We train the model to jointly perform future image generation and action planning  conditioned on past observations and language instructions. 
This task provides a dense and pixel-level supervision signal, compelling the model to learn the causal relationships among instructions, actions, and visual predictions. 
The joint modeling enforces consistency between predictions, enabling mutual supervision and refinement. 
Our model adopts a mixed autoregressive-diffusion transformer architecture~\cite{MoT, janusflow, llamafusion, bagel} to achieve unified \textbf{vision}-\textbf{language} understanding, \textbf{world} modeling, and \textbf{action} planning. 
Specifically, we use the autoregressive pipeline for visual and instruction understanding, and the diffusion pipeline~\cite{rectified-flow, flow-matching} for image and action generation.
We use joint attention to enable interactions across all modalities and employ a Mixture-of-Transformers (MoT) design~\cite{MoT} to effectively decouple the parameters associated with different modalities and enhance the model capacity for joint generation and planning. 
Extensive experiments on the NAVSIM~\cite{navsim-v1, navsim-v2} benchmark show that our model not only achieves superior planning performance but also demonstrates a remarkable ability to generate high-fidelity and instruction-compliant future images and plausible trajectories.
